\documentclass[11pt]{article}
\usepackage[a4paper,margin=25mm]{geometry}
\usepackage{amsmath,amssymb,amsthm,booktabs,array,longtable}
\usepackage[T1]{fontenc}
\usepackage{lmodern}
\usepackage[colorlinks=true,urlcolor=blue,citecolor=blue,linkcolor=blue]{hyperref}
\newtheorem{theorem}{Theorem}
\newtheorem{corollary}[theorem]{Corollary}
\newcommand{\coef}[2]{[#1^{#2}]}
\title{Proving Werner Schulte's coefficient and lattice-path conjectures by integer-parameter extensions of Kruchinin's formula}
\author{Haobo Ma\\\small ChronoAI Pte Ltd, Singapore
\and Wenlin Zhang\\\small National University of Singapore\\\small ChronoAI Pte Ltd, Singapore}
\date{16 November 2026}
\begin{document}
\maketitle
\begin{abstract}
Extending Vladimir Kruchinin's coefficient-transfer formula (2011,
Theorem 8.1), we give a Lagrange-inversion proof of Werner Schulte's
power-coefficient conjecture (OEIS A262082, revision 23, FORMULA,
line 44, beginning ``B(m,n,k) = m/(n*k+m)'', posed 10 September 2015)
for all positive integers $n,m$ and every $k\geq0$.
The alternating-sum conversion is the quadrinomial identity proved by
Roberto Blanco G\'omez (12 October 2026, Section 2).
The coefficient formula extends to all integer parameters, with an explicit
value when $nk+m=0$. An ordered-forest bijection proves Schulte's two
nonnegative-path conjectures. The same argument gives the trinomial
analogues, logarithms and substitutions, and Hanna's binomial-square
exponential identities in the cross-referenced entries. The general
positive-power transfer precedes the conjecture; the trinomial excursion
family also has a later proof by Yan, who identifies the same family
in Yang--Zhang--Yang's earlier work.
\end{abstract}

\section{Power coefficients, including the zero denominator}
Schulte asks for the coefficients of powers of the integer formal series
$P=1+xP^n+x^2P^{2n}+x^3P^{3n}$ in A262082 \cite{schulte}.
We prove the proposed formula and extend it to integer $n$ and integer
powers of $P$. This also gives a precise continuation of A261440's warning
about negative parameters and division by zero \cite{trinomial}.

The general positive-power transfer was already proved by Vladimir
Kruchinin in 2011 \cite[Theorem 8.1, pp.~13--14]{kruchinin1} and retained
in the two-author 2013 revision \cite[Theorem 28, equation (26),
pp.~23--24]{kruchinin2}. In their composita notation,
$B_x^\Delta(u,v)=[x^u](xB(x))^v$. Set $(u,v)=(k+m,m)$ and take
their functional-equation exponent to be $n$. Their auxiliary indices
become $i_{n-1}=nk+m=N$ and $i_n=N+k$; hence their expression is
$m[t^{N+k}](tQ_d(t))^N/N=m[t^k]Q_d(t)^N/N$.
For $n\geq1,m\geq1$ this proves the implicit-series step for every $d$,
not just three fixed rows. The 2011 text also discusses negative
functional-equation exponents via reciprocals; the signed powers and
zero-denominator value below are stated directly with their full domains.
Scheuer's 23 June 2015 answer already pointed to the general theorem
\cite{scheuer}; Kruchinin and Perminova give a later composita treatment
\cite[Theorem 2.3]{perminova}.

Bala's positive binomial sums for $n=1,2,3$ \cite{bala1,bala2,bala3}
are explicit specializations; the $m=1,n=1$ formula is recorded
by Deutsch in 2003 \cite{bala1}. Blanco G\'omez proves the alternating
quadrinomial polynomial identity \cite[Section 2]{blanco}.
Here the bridge from ordinary polynomial powers to the implicit series is
Lagrange inversion \cite[Theorem 2.1.1, equation (2.1.1)]{gessel}.
The integer-parameter form below keeps the boundary $nk+m=0$ rather than
dividing by it. Ordinary polynomial-coefficient and unrestricted-path
interpretations are also treated by Fahssi \cite[Sections 1--2]{fahssi};
the forest interpretation here imposes nonnegativity. The present
contribution is a common proof of the stated formulas and their signed
continuations, with explicit treatment of the removable division by zero,
rather than a priority claim for the positive-power transfer.

For a nonnegative integer $k$, write $\coef{x}{k}F$ for the coefficient
of $x^k$ in a formal series $F$.
For an integer $d\geq1$ put $Q_d(t)=1+t+\cdots+t^d$.
Integer powers, including negative ones, of a series with constant term
$1$ are defined by multiplication and its formal reciprocal. For $A\in\mathbb Z$
and $b\geq0$, use the polynomial convention
\[
 \binom A0=1,\qquad
 \binom Ab=\frac{A(A-1)\cdots(A-b+1)}{b!}\quad(b\geq1).
\]
In particular, $\binom Ab=0$ when $A\geq0$ and $b>A$.
The indicator $\mathbf1_{d+1\mid k}$ is $1$ when $d+1$ divides $k$
and $0$ otherwise.

\begin{theorem}\label{main}
For every $d\geq1$ and $n\in\mathbb Z$ there is a unique
$P_{d,n}\in1+x\mathbb Z[[x]]$ such that
\begin{equation}\label{native}
 P_{d,n}(x)=Q_d\bigl(xP_{d,n}(x)^n\bigr).
\end{equation}
For $m\in\mathbb Z$, the constant coefficient of $P_{d,n}^m$ is $1$.
For $k\geq1$, writing $N=nk+m$, one has
\begin{equation}\label{denominatorfree}
 \coef{x}{k}P_{d,n}^m
   =\frac{m}{k}\coef{t}{k-1}Q_d'(t)Q_d(t)^{N-1}.
\end{equation}
Equivalently,
\begin{equation}\label{boundary}
 \coef{x}{k}P_{d,n}^m=
 \begin{cases}
 \displaystyle\frac{m}{N}\coef{t}{k}Q_d(t)^N,&N\ne0,\\[6pt]
 \displaystyle\frac{m}{k}\bigl(1-(d+1)\mathbf1_{d+1\mid k}\bigr),&N=0.
 \end{cases}
\end{equation}
For $d=2,3$ and $N\ne0$ the first line is
\begin{equation}\label{alternating}
 \coef{x}{k}P_{d,n}^m
 =\frac{m}{N}\sum_{j=0}^{k}
 (-(d-1))^j\binom Nj\binom{dN-2j}{k-j}.
\end{equation}
For $n\geq0,m\geq1$, equation \eqref{alternating} also holds at $k=0$;
the sum may then stop at $\min(k,N)$ and all its remaining upper indices
are nonnegative.
\end{theorem}

\begin{proof}
The coefficient of degree $r$ in
$1+\sum_{i=1}^d x^iP^{in}$ depends only on the coefficients of $P$
of degrees less than $r$. This holds for negative exponents too: if
$P=1+\sum_{r\geq1}p_rx^r$ and $P^{-1}=1+\sum_{r\geq1}q_rx^r$, then
$q_r=-\sum_{i=1}^r p_iq_{r-i}$. Thus coefficient recursion constructs
exactly one integer solution of \eqref{native}.

Put $V=xP_{d,n}^n$. Then $V=xQ_d(V)^n$ and $P_{d,n}=Q_d(V)$.
Over $\mathbb Q[[x]]$, apply Lagrange inversion to
$R(t)=Q_d(t)^n$ and $\varphi(t)=Q_d(t)^m$. Both are formal series
for every integer $n,m$, because $Q_d(0)=1$. For $k\geq1$ it gives
\[
 \coef{x}{k}\varphi(V)
 =\frac1k\coef{t}{k-1}\varphi'(t)R(t)^k
 =\frac{m}{k}\coef{t}{k-1}Q_d'(t)Q_d(t)^{N-1},
\]
proving \eqref{denominatorfree}. If $N\ne0$, use
$(Q_d^N)'=NQ_d'Q_d^{N-1}$ and
$\coef{t}{k-1}H'=k\coef{t}{k}H$.
If $N=0$, use $Q_d'/Q_d=(\log Q_d)'$ and
\[
 \log Q_d(t)=\log(1-t^{d+1})-\log(1-t),\qquad
 \coef{t}{k}\log Q_d(t)=\frac{1-(d+1)\mathbf1_{d+1\mid k}}{k}.
\]
This proves \eqref{boundary}.

For $d=2,3$ the exact identity
$Q_d(t)=(1+t)^d-(d-1)t(1+t)^{d-2}$ yields, for every integer $N$,
\begin{equation}\label{polynomial}
 Q_d(t)^N=(1+t)^{dN}
 \sum_{j\geq0}\binom Nj\left(\frac{-(d-1)t}{(1+t)^2}\right)^j.
\end{equation}
This is a valid formal binomial expansion: the argument has zero
constant term. Terms $j>k$ cannot affect degree $k$, proving
\eqref{alternating}. When $N\geq0$ the series stops at $j=N$.
At $k=0,n\geq0,m\geq1$, $N=m$ and the coefficient formula equals $1$.
\end{proof}

For $d=3,n,m\geq1$, \eqref{alternating} is precisely Schulte's
conjectured $B(m,n,k)$ formula; $m=1$ gives his $A(n,k)$ formula,
including $n=0$. The same specializations for $d=2$ prove the coefficient
conjectures in A261440. Its negative-parameter comment is not a conjecture.
The extension in Theorem~\ref{main} is not asserted there: for example,
\[
 P_{3,-1}=1+x-x^4+O(x^5),\qquad \coef{x}{1}P_{3,-1}=1.
\]
The latter has $nk+m=0$ and uses the second line of \eqref{boundary}.
Signed coefficients occur; no path-count interpretation is assigned to
negative $n$.

\section{Nonnegative paths and trees}
\begin{corollary}\label{paths}
Let $d,n,m$ be positive integers. The coefficient $\coef{x}{k}P_{d,n}^m$ counts paths
of length $nk+m-1$, from height $m-1$ to height $0$, with steps
$(1,in-1)$ for $0\leq i\leq d$, all of whose heights are nonnegative.
For $m=1$ these are excursions of length $nk$.
\end{corollary}
\begin{proof}
Consider rooted plane trees with allowed outdegrees $0,n,2n,\ldots,dn$.
A plane tree has a specified left-to-right order on the children at each
vertex. Give a vertex of outdegree $in$ weight $x^i$, including weight $1$
for a leaf. Root decomposition gives \eqref{native}; an ordered forest
of $m$ such trees has generating function $P_{d,n}^m$.

A forest of weight $x^k$ has total outdegree $nk$. If it has $v$ vertices,
its number of edges is $v-m=nk$, so $v=nk+m$.
Traverse the vertices in preorder, starting with $m$ pending vertex
positions. Visiting an outdegree-$in$ vertex consumes one position and
inserts its $in$ ordered child positions before the remaining positions.
The pending count changes by $in-1$; it is positive before the final
vertex and zero afterwards. The final vertex is a leaf.
Subtract $1$ from the pending count and omit that final leaf to obtain
the required path.

Conversely, append a final down-step to the path and add $1$ to its
heights. The positive pending counts reconstruct the forest by filling
the next position and inserting its children in order. The appended
step exhausts all positions. The reconstructed outdegrees sum to
$v-m=nk$, so its weight is $x^k$. The constructions are inverse.
When $k=0$ the forest consists of $m$ leaves and the path has $m-1$
down-steps.
\end{proof}

For $d=3$ this proves Conjectures 1 and 2 of A262082; for $d=2$ it
proves the two path conjectures of A261440, whose $m=1,n=1,2$ cases
are already identified there as proved.
The tree and ordered-forest mechanism is classical; see
Sokal \cite[Propositions 3.1--3.2, equations (3.4), (3.6)]{sokal}
and Gessel \cite[Section 4.5]{gessel}. Yan's 2025 proof covers
the entire $d=2,m=1$ family \cite[Section 5, Corollary 5.2(D1)]{yan}
and identifies its path class with the 2024 enumeration of
Yang--Zhang--Yang \cite{yang}:
Yan's odd parameter $2n-1$ gives the steps $-1,n-1,2n-1$ and
length $nk$. The precise correspondence used here is the excursion
case $m=1$; Corollary~\ref{paths} supplies the arbitrary-$m$ forest statement.
For $n=m=1,d=3$, replace an outdegree-$i$ vertex, other than the final
leaf, by $i$ up-steps followed by one down-step. This gives Dyck paths
with no four consecutive up-steps, as recorded by Callan in A036765.
The same sequence counts noncrossing partitions with block sizes at
most $3$: a noncrossing partition has no $a<b<c<e$ with $a,c$ in one
block and $b,e$ in another, and its block containing $1$, if of size $i$,
leaves $i$ independent intervals. Hence its generating function satisfies
$P=1+\sum_{i=1}^3x^iP^i$.

\section{Logarithms, substitutions, and ordinary rows}
The logarithmic coefficient mechanism is also classical:
Banderier--Flajolet \cite[Corollary 2]{banderier} express
excursions through bridges, while Gessel \cite{gessel} gives powers
and logarithms directly by Lagrange inversion (Section~2.2,
equations (2.2.4) and (2.2.9)). The specializations below include
Bala's logarithmic formulas, not new positive-parameter identities.
\begin{corollary}\label{logarithm}
For $d\geq1$, $n\ne0$ and $k\geq1$,
\begin{equation}\label{logformula}
 \coef{x}{k}\log P_{d,n}=\frac1{nk}\coef{t}{k}Q_d(t)^{nk}.
\end{equation}
For $n=0$ this coefficient is
$\bigl(1-(d+1)\mathbf1_{d+1\mid k}\bigr)/k$.
For every pair of integers $n,a$,
\begin{equation}\label{composition}
 P_{d,n}(x)=P_{d,a}\bigl(xP_{d,n}(x)^{n-a}\bigr),\qquad
 P_{d,n}\left(\frac{x}{Q_d(x)^n}\right)=Q_d(x).
\end{equation}
For $d=2$ and every integer $n$, one also has
\begin{equation}\label{rational}
 P_{2,n}=\frac{(1+xP_{2,n}^n)^2}{1+xP_{2,n}^{n-1}}.
\end{equation}
\end{corollary}
\begin{proof}
With $V$ as in the theorem, Lagrange inversion for $\log Q_d$ gives
\[
\coef{x}{k}\log P_{d,n}=k^{-1}\coef{t}{k-1}Q_d'Q_d^{nk-1}.
\]
Differentiating $Q_d^{nk}$ proves \eqref{logformula} when $n\ne0$.
When $n=0$, $P_{d,0}=Q_d$.

Set $Y=xP_{d,n}^{n-a}$. It has linear coefficient $1$, and
$P_{d,n}=Q_d(YP_{d,n}^a)$. Coefficient recursion in $x$ uniquely
determines a unit solution of this equation, so it equals $P_{d,a}(Y)$.
For the second identity put $g(t)=t/Q_d(t)^n$. It has linear coefficient
$1$ and $g(V(x))=x$, hence also $V(g(x))=x$.
Consequently $P_{d,n}(g(x))=Q_d(x)$.
Finally, put $z=xP_{2,n}^n$. Then
$P_{2,n}(1+xP_{2,n}^{n-1})=P_{2,n}+z=(1+z)^2$.
The denominator is a unit, proving \eqref{rational}.
\end{proof}

For $d=2,3$ and $n>0$, the exponential form of \eqref{logformula} is
\begin{equation}\label{correctedexp}
 P_{d,n}(x)=\exp\left(\sum_{k\geq1}\frac{x^k}{nk}
 \sum_{j=0}^k (-(d-1))^j\binom{nk}{j}\binom{dnk-2j}{k-j}\right).
\end{equation}
The exponential lines in A262082 and A261440 omit $x^k$.
Without it the displayed exponent is not a formal series of positive
order; \eqref{correctedexp} supplies the corrected statement.
Equation \eqref{composition} also gives Hanna's reversion formulas in
A186241 and A200731: $xP_{3,n}^n$ is the inverse of $x/Q_3(x)^n$.
For $n\geq1$, putting $U=xP_{d,n}(x^n)$ gives
$U=xQ_d(U^n)$, so $U$ is the inverse of $x/Q_d(x^n)$ as well.
At $d=3,n=2$ this is Bala's 27 July 2023 reversion formula in A186241.

\begin{corollary}\label{rows}
For $N\geq0$ put $C_{N,k}=\coef{t}{k}Q_3(t)^N$, with
$C_{N,k}=0$ for $k<0$. Then
\begin{align}
 C_{N,k}&=\sum_{b=0}^{\lfloor k/2\rfloor}\binom Nb\binom N{k-2b},\label{positive}\\
 C_{N,k}&=\sum_{i=0}^{\lfloor k/4\rfloor}(-1)^i\binom Ni
                 \binom{N+k-4i-1}{k-4i}\quad(N\geq1).\label{inclusion}
\end{align}
Their support is $0\leq k\leq3N$, with positive values throughout;
$C_{N,0}=C_{N,3N}=1$ and $C_{N,k}=C_{N,3N-k}$ on that support.
They satisfy
\[
 C_{N+1,k}=\sum_{i=0}^3 C_{N,k-i},\qquad
 \sum_k C_{N,k}=4^N,\qquad \sum_k(-1)^k C_{N,k}=0\quad(N>0).
\]
For $n\geq0,m\geq1,N=nk+m$, the power coefficient is
$mC_{N,k}/N$, so \eqref{positive} gives Bala's formulas at $n=1,2,3$.
\end{corollary}
\begin{proof}
Use $Q_3(t)=(1+t)(1+t^2)$ for \eqref{positive} and
$Q_3(t)=(1-t^4)/(1-t)$ for \eqref{inclusion}.
Choosing one term from each of $N$ factors also identifies $C_{N,k}$
with compositions into $N$ parts in $\{0,1,2,3\}$, so precisely that
support is positive. Replacing each part $i$ by $3-i$ gives symmetry.
Multiplication by $Q_3$, then evaluation at $1$ and $-1$, gives the
recurrence and row sums. The last assertion is Theorem~\ref{main}.
\end{proof}
These ordinary-row identities are recorded in A008287 \cite{quadrinomial},
including Shevelev's 2010 formulas, Bala's 2013 inclusion--exclusion sum,
and the composition, unrestricted-path and four-sided-die interpretations.
Myerson's 2014 coefficient-extraction argument \cite{myerson} gives
the same bounded-composition inclusion--exclusion mechanism; the
summation conventions are discussed further by Belbachir--Mendon\c ca
\cite[Section 3.3, Remark 3.2 and Proposition 3.3]{belbachir}.
The alternating identity \eqref{polynomial} at $d=3,N\geq0$ is reproved
here with credit to Blanco G\'omez's 2026 proof.

\section{Hanna's binomial-square identities}
The cross-referenced entries contain two further families of exponential
identities. Define polynomials
\[
 S_r(v)=\sum_{j=0}^r\binom rj^2v^j,\qquad
 R_r(v)=\sum_{j=0}^{r-1}\binom{r-1}j\binom rj v^j\quad(r\geq1).
\]
\begin{corollary}\label{hanna}
For every integer $n$, with $P=P_{3,n}$,
\begin{equation}\label{hanna1}
 \log P=\sum_{r\geq1}\frac{(xP^{n-1})^r}{r}S_r(xP^n).
\end{equation}
For $P=P_{3,1}$ one also has
\begin{equation}\label{hanna2}
 \log\bigl((1-x)P\bigr)
 =\sum_{r\geq1}\frac1r\left(\frac{x^2P}{(1-x)^2}\right)^rR_r(x).
\end{equation}
The infinite-binomial-sum versions follow from
\begin{equation}\label{eulertransform}
 \sum_{j\geq0}\binom{a+j}j\binom{b+j}jv^j
 =\frac{\sum_j\binom aj\binom bjv^j}{(1-v)^{a+b+1}}
 \qquad(a,b\geq0).
\end{equation}
\end{corollary}
\begin{proof}
Over $\mathbb Q[[v]][[u]]$, let $T=u(1+T)(1+vT)$ and
$Y=T/u=(1+T)(1+vT)$. Lagrange inversion gives
\[
 \coef{u}{r}\log Y=\frac1r\coef{t}{r}(1+t)^r(1+vt)^r
 =\frac{S_r(v)}r,
\]
and
\[
 \coef{u}{r}\log(1+T)
 =\frac1r\coef{t}{r-1}(1+t)^{r-1}(1+vt)^r
 =\frac{R_r(v)}r.
\]
In the first equality for $\log Y$, apply the derivative form to
$\log((1+t)(1+vt))$ and then differentiate the $r$th power.
For \eqref{hanna1}, substitute $u=xP^{n-1}$, $v=xP^n$, and $T=xP^n$;
then $Y=P$. All substitutions have positive order.
For \eqref{hanna2}, put $Y=(1-x)P$, $T=Y-1$, $v=x$ and
$u=x^2P/(1-x)^2$. The native equation gives
\[
 T=x^2P^2(1+xP)=u(1+T)(1+xT).
\]
Thus the second coefficient identity applies.

For completeness, the left side of \eqref{eulertransform} equals
$\frac1{b!}(\frac{d}{dv})^b[v^b(1-v)^{-a-1}]$.
Leibniz's rule, followed by multiplication by $(1-v)^{a+b+1}$, gives
\[
 \sum_{i=0}^b\binom bi\binom{a+i}i v^i(1-v)^{b-i}.
\]
For $0\leq j\leq b$ its coefficient is
$\binom bj\sum_{i=0}^j(-1)^{j-i}\binom ji\binom{a+i}a
=\binom bj\binom aj$; the last equality follows by $j$ successive
applications of Pascal's identity. This proves \eqref{eulertransform}.
Take $a=b=r$ for $S_r$ and $a=r-1,b=r$ for $R_r$.
\end{proof}
At $n=1,2,3$, \eqref{hanna1} and its transformed version reproduce
Hanna's formulas in A036765, A186241 and A200731; the first two
entries date the relevant contributions to February and November 2011.
The individual undated formulas in A200731 are not assigned an insertion
date here. Equation \eqref{hanna2} gives the second
pair in A036765, dated 24 February 2011.
More generally the same bivariate calculation proves Hanna's $p,q$
comment in A200731: for integer $p,q$, a unit series $F$ satisfies
$\log F=\sum_{r\geq1}(xF^p)^rS_r(xF^q)/r$ if and only if
$F=(1+xF^{p+1})(1+x^2F^{p+q+1})$.
Indeed $Y$ is the unique unit solution of
$Y=(1+uY)(1+uvY)$; substitute $u=xF^p,v=xF^q$ in this identity and
in its logarithm. Uniqueness and injectivity of the formal logarithm
give both implications.

\section{Source coverage and verification scope}
The following locators refer to the revisions specified in the references.
``Proved'' describes an assertion actually made by the source; the
integer-parameter continuation is an extension, not a refutation of a
statement made only for positive parameters.
At A262082 line 44 the power formula is explicitly labelled
``(conjectured)''; its COMMENTS lines 22--31 name ``Conjecture 1''
and ``Conjecture 2'' and specify the lengths, steps and nonnegativity
used in Corollary~\ref{paths}. A261440 line 43 labels its power formula
``(conjectured)'', and COMMENTS lines 23--33 pose the analogous path
statements. These labels describe the posing text, not the absence of
earlier encompassing theorems or later proofs.
\setlength{\LTpre}{3pt}
\begin{longtable}{@{}p{0.43\linewidth}p{0.52\linewidth}@{}}
\toprule Source item & Coverage\\\midrule\endhead
A262082, FORMULA lines 32, 44: $A(n,k)$; $B(m,n,k)$ & Proved, Theorem~\ref{main}, including $k=0$ and the $n=0$ row. Prior general transfer: Kruchinin, 2011; revised 2013; later treatment Kruchinin--Perminova, 2021.\\
A262082, COMMENTS: Conjectures 1, 2 & Proved, Corollary~\ref{paths}.\\
A262082, FORMULA: substitutions and inverse substitution & Proved, Corollary~\ref{logarithm}.\\
A262082, FORMULA: exponential without $x^k$ & Corrected formal-series statement, equation~\eqref{correctedexp}.\\
A262082, FORMULA: degrees $0$--$5$ & Proved by coefficient extraction in Theorem~\ref{main}.\\
A261440, FORMULA and COMMENTS: coefficients, two path conjectures, rational form and substitutions & Proved, $d=2$ in Theorem~\ref{main} and Corollaries~\ref{paths}, \ref{logarithm}. Prior proof cited for the full $m=1$ excursion family: Yan, 2025, relating it to Yang--Zhang--Yang, 2024; $m=1,n=1,2$ were already noted as proved by the source.\\
A261440: exponential; negative-$n$ comment & Corrected exponential as above. Integer parameters and $nk+m=0$: not posed; extension proved in Theorem~\ref{main}.\\
A036765, A186241, A200731: power sums and logarithms & Reproved; prior formulas Bala, 2015. The $m=n=1$ coefficient formula is recorded by Deutsch, 2003.\\
A036765, A186241, A200731: reversion and binomial-square exponential formulas; A200731 $p,q$ comment & Reproved; prior formulas of Arndt and Hanna (dated 2011 where specified), and Bala, 2023; Corollaries~\ref{logarithm}, \ref{hanna}.\\
A036765: bounded-degree trees, Dyck paths and noncrossing partitions & Proved as consequences of Corollary~\ref{paths}, agreeing with the interpretations recorded there.\\
A008287, FORMULA: alternating sum & Reproved; prior proof Blanco G\'omez, 2026, Section 2.\\
A008287: positive sum, inclusion--exclusion, symmetry, recurrence, row sums and elementary counting interpretations & Reproved, Corollary~\ref{rows}; prior formulas Shevelev, 2010, and Bala, 2013. Related bounded-composition proofs: Myerson, 2014; Belbachir--Mendon\c ca, 2026.\\
Other cross-referenced recurrences, asymptotics, production matrices, transforms and avoidance models & Not covered by these arguments; no new openness claim.\\
\bottomrule
\end{longtable}

\noindent\textbf{Formal scope.}
The pinned file
\href{https://github.com/the-omega-institute/trureturing/blob/2f8421f6b90d37b4815b887f6c60987553831b20/D5/S3/Combinatorics/PowerSeries/SchulteImplicitQuadrinomialPowers.lean}{\texttt{SchulteImplicitQuadrinomialPowers.lean}}
defines \texttt{Native} by \eqref{native} for $d=3$.
Its declaration \texttt{result} proves integer-series unique existence
and \eqref{alternating} for $n,m\geq1$ and every $k\geq0$, with the
integer coefficient mapped to $\mathbb Q$. Binomial indices in that
declaration use natural subtraction; terms $j>N$ vanish. Since
$k\leq3N$ and $j\leq\min(k,N)$, the nonzero terms agree with the
ordinary-binomial convention here.

The pinned file
\href{https://github.com/the-omega-institute/trureturing/blob/2f8421f6b90d37b4815b887f6c60987553831b20/D5/S1/Recurrence/SchulteQuadrinomialAlternatingBinomialExpansion.lean}{\texttt{SchulteQuadrinomialAlternatingBinomialExpansion.lean}}
defines \texttt{T N k} as the degree-$k$ integer polynomial coefficient
of $Q_3^N$. Its declaration \texttt{result} proves the alternating
formula for $N\geq0,0\leq k\leq3N$.
The integer-parameter extension, zero-denominator value, paths,
trinomial identities, substitutions and binomial-square identities have
the written proofs above. The accompanying \texttt{verify.py} compares
coefficient recursion, integer reciprocals, polynomial powers, path
counts, logarithms and compositions using exact arithmetic;
\texttt{verification.json} specifies the finite ranges.

\begin{thebibliography}{19}
\bibitem{kruchinin1} V. Kruchinin, \emph{Compositae and their properties},
arXiv:1103.2582v1 (14 March 2011), Section 8, Theorem 8.1, pp.~13--14.
\url{https://arxiv.org/abs/1103.2582v1}.
\bibitem{kruchinin2} V. V. Kruchinin and D. V. Kruchinin,
\emph{Composita and its properties}, arXiv:1103.2582v2 (25 March 2013),
Section 6, Theorem 28, equation (26), pp.~23--24.
\url{https://arxiv.org/abs/1103.2582v2}.
\bibitem{scheuer} M. Scheuer, answer to \emph{Asymptotics of coefficients
of implicitely defined generating function}, MathOverflow,
23 June 2015, final displayed equation and preceding theorem pointer.
\url{https://mathoverflow.net/a/209984}.
\bibitem{perminova} D. V. Kruchinin and M. Y. Perminova,
\emph{About solving some functional equations related to the Lagrange
inversion theorem}, Montes Taurus J. Pure Appl. Math. \textbf{3}(1)
(2021), 62--69, Theorem 2.3, pp.~65--66.
\url{https://mtjpamjournal.com/papers/article_id_mtjpam-d-20-00011}.
\bibitem{sokal} A. D. Sokal, \emph{A ridiculously simple and explicit
implicit function theorem}, arXiv:0902.0069v1 (31 January 2009),
Propositions 3.1--3.2 and equations (3.4), (3.6), pp.~10--12;
S\'eminaire Lotharingien de Combinatoire \textbf{61A}, article B61Ad.
\url{https://arxiv.org/abs/0902.0069v1}.
\bibitem{banderier} C. Banderier and P. Flajolet,
\emph{Basic analytic combinatorics of directed lattice paths},
Theoret. Comput. Sci. \textbf{281} (2002), 37--80, Corollary 2
and Example 3, equation (24).
\url{https://doi.org/10.1016/S0304-3975(02)00007-5}.
\bibitem{yang} L. Yang, Y.-Y. Zhang and S.-L. Yang,
\emph{Enumeration of the Motzkin paths above a line of rational slope},
Discrete Math. \textbf{347}(7) (2024), 114013.
The identification of its path class used here is given in
\cite[Section 5]{yan}.
\url{https://doi.org/10.1016/j.disc.2024.114013}.
\bibitem{yan} J. Yan, \emph{Lattice paths enumerations weighted by
ascent lengths}, arXiv:2501.01152v1 (2 January 2025), Theorem 3.1,
pp.~5--6, and Section 5, Corollary 5.2(D1), pp.~12--13;
Discrete Math. \textbf{349}(2) (2026), 114783.
\url{https://arxiv.org/abs/2501.01152v1}.
\bibitem{myerson} G. Myerson, answer on bounded positive parts,
Math.StackExchange, 6 July 2014, generating-function factorization
and coefficient extraction.
\url{https://math.stackexchange.com/a/857682}.
\bibitem{belbachir} H. Belbachir and J. R. G. Mendon\c ca,
\emph{Bounded compositions and products of chains, or bi(s)nomials
everywhere from de Moivre's formula to spin multiplicities and
constant-sum codes}, arXiv:2610.31865v1 (22 October 2026),
Section 3.3, equations (13)--(17), Remark 3.2, Proposition 3.3
and Table 2, pp.~4--6.
\url{https://arxiv.org/abs/2610.31865v1}.
\bibitem{schulte} W. Schulte, \emph{OEIS A262082}, 10 September 2015,
revision 23 (3 June 2018), COMMENTS and FORMULA, especially
\href{https://github.com/oeis/oeisdata/blob/49db7744e1fdbefbe376d1d018370c017d4c288b/seq/A262/A262082.seq#L44}{the power-coefficient formula, line 44}.
\url{https://oeis.org/A262082}.
\bibitem{trinomial} W. Schulte, \emph{OEIS A261440}, 18 August 2015,
revision 32 (3 June 2018), COMMENTS and FORMULA.
\url{https://oeis.org/A261440}.
\bibitem{bala1} \emph{OEIS A036765}, revision 195 (30 October 2026),
COMMENTS and FORMULA: Deutsch (29 November 2003), Hanna (13 and
24 February 2011), Arndt (10 June 2011), Bala (21 June 2015).
\url{https://oeis.org/A036765}.
\bibitem{bala2} \emph{OEIS A186241}, revision 58 (11 July 2026),
FORMULA: Hanna (11 November 2011), Bala (21 June 2015; 27 July 2023).
\url{https://oeis.org/A186241}.
\bibitem{bala3} \emph{OEIS A200731}, revision 19 (2 November 2019),
COMMENTS and FORMULA: Hanna (entry dated 21 November 2011),
Bala (21 June 2015).
\url{https://oeis.org/A200731}.
\bibitem{gessel} I. M. Gessel, \emph{Lagrange inversion},
J. Combin. Theory Ser. A \textbf{144} (2016), 212--249;
arXiv:1609.05988v1 (20 September 2016), Theorem 2.1.1.
\url{https://arxiv.org/abs/1609.05988v1}.
\bibitem{quadrinomial} \emph{OEIS A008287}, revision 172
(30 October 2026), COMMENTS and FORMULA.
\url{https://oeis.org/A008287}.
\bibitem{blanco} R. Blanco G\'omez,
\emph{A binomial-theorem proof of Schulte's identity for quadrinomial
coefficients (OEIS A008287)}, version 1.0, 12 October 2026, Section 2.
\url{https://doi.org/10.5281/zenodo.23315738}.
\bibitem{fahssi} N.-E. Fahssi, \emph{Polynomial Triangles Revisited},
arXiv:1202.0228v7 (25 July 2012), Sections 1--2.
\url{https://arxiv.org/abs/1202.0228v7}.
\end{thebibliography}
\noindent\small Copyright 2026 Haobo Ma and Wenlin Zhang.
Licensed under Creative Commons Attribution 4.0 International (CC BY 4.0).
\url{https://creativecommons.org/licenses/by/4.0/}.
\end{document}
